% \section*{Lernziele Selbstkorrigierende Codes}
% \begin{todolist}
% \item Ich weiss, wie die Prüfziffer von EAN-Codes definiert ist.
% \item Ich kann die Prüfziffern eines EAN-Codes berechnen.
% \item Ich kann entscheiden, ob ein gegebener EAN-Code korrekt ist.
% \item Ich kann bestimmen, welche Fehler (Vertippen einer Ziffer, Vertauschen von zwei Ziffern, Ziffer vergessen) durch eine Prüfziffer erkannt werden.
% \item Ich kann den Abstand zweier Bitfolgen gleicher Länge bestimmen.
% \item Ich kann den Abstand einer Kodierung bestimmen.
% \item Ich kann bestimmen, wie viele Fehler eine Kodierung erkennt.
% \item Ich kann bestimmen, wie viele Fehler eine Kodierung korrigieren kann.
% \item Ich kann erklären, welchen Abstand \textit{k}-Fehlererkennende Kodierungen haben müssen.
% \item Ich kann erklären, welchen Abstand \textit{k}-Fehlererkorrigierende Kodierungen haben müssen.
% \item Ich kann den \textit{n}-dimensionalen Hyperwürfel für kleine \textit{n} zeichnen.
% \item Ich kann Kodierungen angeben, die einen bestimmten Abstand haben, bzw. eine bestimmte Anzahl Fehler erkennen oder korrigieren können.
% \item Ich kann die ``Kartentrick-Kodierung'' einer Bitfolge angeben.
% \item Ich kann erklären, wie die optimale Kartentrickkodierung ( = optimale Länge und Breite für das Rechteck) für eine Bitfolge mit einer gegebenen Länge aussieht.
% \item Ich kann überprüfen, ob eine gegebene Kartentrick-Kodierung korrekt ist.
% \item Ich kann bei einer Kartentrick-Kodierung mit einem Fehler den Fehler finden und korrigieren.
% \end{todolist}



\section{Elektronische Schaltungen und Computerhardware}
Stoffumfang: Notizen 0, Notizen 1 und Notizen 2 (Moodle) sowie das im Unterricht besprochene.
Die Demonstration der Addition auf dem 8-Bit Computer ist nicht Prüfungsstoff.
\begin{todolist}
\item Ich kenne den konzeptionellen Aufbau eines Transistors.
\item Ich weiss, wie ein Transistor als Schalter ohne bewegliche Teile verwendet werden kann.
\item Ich kann erklären, warum Transistoren so eine wichtige Rolle in moderner Hardware spielen.
\item Ich kenne den elektronischen Aufbau, die Logik (Wahrheitstabellen) und die Schaltsymbole der:
  \begin{todolist}
  \item Inverter
  \item UND-Gatter
  \item ODER-Gatter
  \item XOR-Gatter
  \end{todolist}
\item Ich kann erklären, wie ein Computer Informationen abspeichern kann (Memory).
\item Ich kann erklären, wozu ein Computer Speicher (Memory) benötigt.
\item Ich weiss im Detail, wie ein Addierer eines Computers aufgebaut ist und arbeitet (Halbaddierer, Volladdierer).
\item Ich verstehe, wie ein $n$-Bit Addierer aufgebaut ist.
\end{todolist}

\section*{Lernziele Selbstkorrigierende Codes}
Stoffumfang: Seiten 100 bis und mit 106 sowie die Zusammenfassung auf Seite 107 (oben) im  \enquote{Data Science und Sicherheit}
\begin{todolist}
\item Ich weiss, wie EAN-Codes definiert sind und konstruiert werden.
\item Ich weiss, wie ISBN-Codes definiert sind und konstruiert werden.
\item Ich kann bestimmen, welche Fehler (Vertippen einer Ziffer, Vertauschen von zwei Ziffern, Ziffer vergessen) durch eine Prüfziffer erkannt werden.
\item Ich verstehe die Musterlösungen der Aufgaben im Buch und beherrsche die dort gezeigten Rechentechniken und Schritte (modulares Rechnen, Teilbarkeit, Quersummen, \ldots).
\end{todolist}
